\documentclass[10pt,letterpaper]{article}
\usepackage{technical-testing}
\setDocTitle{Animation Benchmark And Regression Procedure}
\setDocSubtitle{Controlled execution and review procedure}
\setDocVersion{1.0}
\setDocStatus{Educational / Draft}
\setDocOwner{Game Development}
\setDocDate{2026-10-30}
\setDocSummary{Controlled execution and review procedure}
\begin{document}
\makedoctitle
\makedocsummary
\tableofcontents\clearpage

\section{Experimental control}
Choose CPU model, GPU, OS, driver, engine version, scene revision, rig count, render and simulation rates, instrumentation overhead and capture mode. Warm caches consistently; state thermal throttling and scene loading conditions.
\section{Metrics}
Animation state update duration (ms), sample/blend duration (ms), p50/p95/p99 evaluation, worst frame spike, total animated characters, asset memory and streaming stalls, foot drift (world units), visually detectable marker timing discrepancy (frames/ms), and game-input-to-cue latency (ms).
\section{Procedure}
\begin{enumerate}
\item Validate functional correctness before performance comparisons. Freeze benchmark fixture and deterministic seed.
\item Perform at least one warm-up run; capture a stated number of measured updates and report repetitions.
\item Compare baseline and treatment under identical hardware/load. Export raw samples and aggregation script.
\item Report confidence interval or distribution spread when supported; never invent results or use unmeasured thresholds.
\item Perform a visual audit at the intended device display refresh and review contact sheet at fractional speed.
\end{enumerate}
\section{Regression gates}
Require both numeric threshold and qualitative acceptance: no state deadlocks, no cue loss, no visible contact skating beyond agreed tolerance, and no unacceptable input latency. All failures link to reproducible traces and source revisions.
\end{document}
